\documentclass[11pt,a4paper]{article}
\usepackage[utf8]{inputenc}
\usepackage[T1]{fontenc}
\usepackage[spanish,es-noshorthands]{babel}
\usepackage{lmodern}
\usepackage{microtype}
\usepackage[margin=2.7cm]{geometry}
\usepackage{parskip}

% Preguntas del entrevistador en negrita. Cambia aquí si prefieres otro formato.
\newcommand{\preg}[1]{\textbf{#1}}

\begin{document}

\begin{center}
{\LARGE\bfseries Gustavo Ariel Schwartz: «Puse a Maradona y al papa y en medio aparecieron los últimos tres presidentes argentinos»\par}
\vspace{1em}
{\large Por Juan José Gómez Cadenas\par}
\end{center}
\vspace{1.5em}

El castellano reserva la palabra «mestizo» para lo que nace de mezclar, y a Gustavo Ariel Schwartz le sienta como un guante. Este argentino que se declara extranjero en todas partes, y que paga ese precio contento, es a la vez un físico (investigador del CSIC en el Centro de Física de Materiales de San Sebastián e investigador asociado al Donostia International Physics Center) y un escritor que lleva quince años al frente de Mestizajes, el programa del DIPC desde el que sienta a la misma mesa a científicos, escritores, filósofos y artistas.

Una de las facetas más destacables de su personalidad (yo diría que extraordinaria, en el mundo de las letras y las ciencias, donde todos andamos bien servidos de ego) es una humildad no impostada que se entiende mejor después de leer \emph{Conexiones}, su último libro, donde sostiene que la creatividad individual es un invento del Romanticismo. De esa idea, de \emph{La entrevista}, la obra de teatro que escribió a cuatro manos con Luisa Etxenike, de Borges y Cortázar, del azar y de las redes con las que ahora cartografía la cultura hablamos durante más de hora y media, hasta llegar a una pregunta que no estaba en el guion, la de qué tienen que ver Messi y Dios.

\preg{Unos versos de Machado dicen «Castilla miserable, ayer dominadora, / envuelta en sus andrajos desprecia cuanto ignora». Son versos muy duros, muy amargos. ¿Describen todavía la España que tú conoces? ¿Describen de alguna manera a Euskadi? ¿Y a la Argentina?}

Diría que quizás describen un poco a la humanidad. Despreciar lo que uno ignora es casi la marca de la casa, y los tiempos que estamos viviendo lo muestran.

\preg{Y al revés, la ignorancia casi implica el desprecio.}

Sí, pero te voy a decir más. Ignorar significa desconocer algo, y sobre esa ignorancia yo puedo decir «desconozco ese algo, me interesa conocerlo y hago un esfuerzo por conocerlo», o puedo simplemente despreciar esa oportunidad de conocer lo desconocido. Y eso es lo que me parece más terrible, no tanto la ignorancia en sí como la falta de ganas, llamémosle, de superar esa ignorancia.

\preg{Todo tu trabajo de décadas parece una pelea desesperada contra esa ignorancia irredenta.}

Sí, es cierto. Esas ganas de conocer son básicamente lo que me llevó a hacer casi todo lo que hago. He tenido y tengo todavía, y cada vez más, esa curiosidad. Cuanto más conocemos, más conscientes somos de todo lo que ignoramos, y en mi caso particular eso me lleva a tener cada vez más ganas de conocer.

\preg{Nadie es profeta en su tierra, dice el refrán. ¿Cuál es tu tierra, Gustavo? Porque a lo mejor te pasa como al que te entrevista, que vives en una función de onda que no se quiere colapsar entre patrias chicas.}

Mirá, lo hablábamos el otro día con Silvina y con Julieta, con mi mujer y con mi hija. Decíamos que nosotros tres tenemos que asumir que somos extranjeros en todos lados, y casi que prefiero eso a ser de un sitio. Hace muchos años me invitaron al Palacio de Aiete a dar una charla que titulé «Reflexiones sobre la identidad». Ahí traía a colación un viejo chiste de Mafalda, que se preguntaba si uno es de donde nace o de donde elige ser. Es un poco esa idea de que la identidad se construye, y de que los que tuvimos la suerte de vivir en muchos sitios nos hemos construido una especie de identidad mestiza. Es una identidad construida con retazos de cada uno de esos sitios donde has vivido. Te iba a decir que te llevas lo mejor de cada sitio, pero al final te llevas lo que eliges, y uno trata de elegir lo mejor. Perteneces en parte a todas partes, porque tienes este pedazo de aquí y este pedazo de allá, y en parte no.

\preg{Y eres forastero en todas partes. Es un extraño precio a pagar.}

Sí. Tiene un precio, pero yo lo pago contento.

\preg{Yo siento lo mismo sin haberme movido tanto. Y Mafalda, que debería ser obligatoria en un doctorado de Filosofía, empalma muy bien con Max Aub, que decía que uno es de donde hizo el bachillerato, que es una forma de negar la patria grande a cambio de la patria chica. Y luego está aquello, todavía más radical, de que uno es de donde cuelga su sombrero. Me parece importante porque vivimos en un tiempo de identitarismos, de nacionalidades, de trincheras.}

De muros, de fronteras.

\preg{De fronteras dentro de fronteras. Quizás está poco de moda no querer estar en ellas, pero es lo que hay. Vamos al \emph{Martín Fierro}. «Cruz no consiente / que se cometa el delito / de matar ansí un valiente». Me quedé asombrado cuando lo leí. Es un policía que tiene acorralado, junto a otros, a un forajido que huye, y no le parece una pena ni le parece mal, le parece un delito matar a un valiente, y se cambia de bando. Pone el valor personal por delante de casi todo. Ese valor a mí me parece que se lee hoy como el de atreverse a decir lo que nadie más se atreve a decir, el de ir a contracorriente, el de no tener una cuenta de TikTok.}

Me parece bien el valor, me parece bien ir a contracorriente; no me parece bien el valor individual, y sobre todo viniendo del \emph{Martín Fierro}. Una crítica que tengo con la Argentina es que, igual que muchos países que fueron colonia, está muy necesitada del caudillismo y es una cosa con la que yo no puedo acordar, porque para mí las soluciones de una sociedad tienen que ser necesariamente colectivas. Entonces el valor vale, valga la redundancia, pero me parece que solamente con ese valor individual no se puede construir una sociedad.

\preg{Y sin embargo es necesario. Sin el individuo que se atreve a decir «yo por aquí no paso», la sociedad tampoco avanza.}

Sin duda, y yo creo que lo que mueve el mundo al final es ese valor individual, pero eso necesita un contexto social que le permita gestionarlo y, si querés, de alguna manera repartirlo socialmente. Volviendo a la Argentina, yo creo que a nivel individual tiene gente extremadamente valiosa, pero es muy difícil que esas individualidades funcionen de forma colectiva, de forma coordinada.

\preg{Es una de las grandes tragedias que uno ve en Argentina. ¿Cómo es posible que un país lleno de gente extraordinaria, cultísima, fantástica, no vaya mejor? Y eso me lleva a otro argentino. Borges, hablando de España, escribió sobre la hombría de bien y el «inútil coraje». Cuando lo leí pensé que estaba hablando de mí, y también de ti, porque todo lo que te he visto hacer desde que te conozco, los nodos, las redes, la conversación, ese intentar constantemente cosas profundas, necesarias y difíciles, es algo de lo que el mundo no quiere saber nada. El mundo está en su casa haciendo un reel de TikTok con un gato. Coraje inútil donde los haya.}

Mirá, en alguno de los proyectos de Mestizajes surgió una frase, a ver si me la acuerdo de memoria, que decía que la condición humana consiste en un creer inútil. Yo creo que cualquier gran idea al principio fue ridícula o inútil, después era peligrosa y finalmente era obvia. Todas las grandes ideas han pasado por esas tres etapas, y es que si no, no puede ser una gran idea. Si una idea desde el primer momento parece útil, necesaria, y todos estamos de acuerdo, entonces no es una gran idea.

\preg{Admitirás que atreverse a empujarla implica dosis muy considerables de esa valentía, porque remas contra un entorno hostil que va a ignorarte o a tirarte piedras. A mí me gusta la frase de Borges porque conecta ese coraje con la hombría de bien, que no hace falta leer en sentido masculino. Uno lo hace porque cree que debe hacerlo, porque hay una necesidad estética y ética.}

O simplemente porque no puede evitarlo. Yo creo que es más eso, que es más íntimo. Hay algo interno que te mueve a hacerlo. Es una cosa que traté de inculcarle mucho a mi hija, que la motivación siempre tiene que ser interna. No puede ser externa, no puede ser el dinero o los premios. Todo eso, si llega, genial, pero no puede ser esa la motivación principal. La motivación principal o te sale de dentro o no vas a ir muy lejos. Yo siempre creo que el motor es interno. Uno no lo puede pensar como que esto que estoy haciendo es necesario para la humanidad. Sería demasiado egocéntrico. Puede ser necesario para mí y después, si eso le sirve a alguien más, estupendo. De eso se habla en el libro con una imagen que a mí me encanta y que pongo en dos lugares diferentes: uno de los personajes verifica una teoría de Aristóteles y duerme tranquilo; después está Galileo, que la contradice y no puede dormir. Yo creo que es eso.

\preg{Es reducir algo tan potente a la frase de Jessica Rabbit. Lo hago porque me dibujaron así.}

Sí, yo creo que al final es poco lo que elegimos. Hay un margen ahí entre lo que viene dado y el libre albedrío, llamémosle, pero es un margen relativamente estrecho.

\preg{Sigo con Borges. Siendo el siglo \textsc{xx} tan rico en escritores, este personaje se sale del tiesto. ¿Qué opinas tú?}

A ver, yo empecé leyendo a Sábato, después a Borges y después a Cortázar. Los tres me gustaron por motivos diferentes, pero yo siempre lo veo a Borges como el último gran clásico y a Cortázar como el primer escritor moderno. Yo creo que los dos han sido maravillosos, pero en distintos, no sé si llamarles estilos o corrientes históricas.

\preg{No lo había formulado nunca así y me gusta mucho. Borges es Planck, es Einstein incluso…}

El tipo se comió todo, se comió toda la filosofía, se comió toda la ciencia.

\preg{…y Cortázar es Bohr y es Heisenberg.}

Y además, cuando mirás un poco desde el punto de vista filosófico, Cortázar está mucho más cerca de Derrida y Borges está más cerca de los platónicos.

\preg{En «Hombre de la esquina rosada», Francisco Real, el Corralero, llega a desafiar al matón oficial, Rosendo Juárez, que aparentemente se muere de miedo y renuncia a pelear. Treinta y cinco años después Borges reescribe la historia desde el punto de vista del cobarde, que dice que no tuvo miedo, que de repente le pareció absurdo liarse a puñaladas por aquello. El Corralero es literalmente el caudillismo del que hablábamos, y a Juárez de repente se lo come la sensación de absurdo y renuncia a jugar ese juego. Parece que Borges nos está diciendo que hay otra elección posible, un mundo menos basado en el coraje inútil y más en la capacidad de hacer las cosas racionalmente. ¿Cómo evalúas ese juego de espejos?}

La verdad es que no sé muy bien qué decirte, porque al final uno puede interpretarlo como un remake de esa historia. Pasaron treinta y cinco años, él pudo haber cambiado su punto de vista, pudo darle otra vuelta y mirarlo de otra manera, o simplemente cambiaron los tiempos y él necesita adaptarse a eso. Es difícil atribuir intenciones. Cualquier análisis que uno pueda hacer es un poco entrar dentro de la ficción, y atribuir cuál fue la intención real, si él no lo dice, es medio complicado. Yo creo que tiene más que ver con el cambio de época que con otra cosa.

\preg{No te voy a preguntar por tu mujer, a quien conozco y admiro, pero sí quiero darle la vuelta metafórica a la Lujanera. Tú le has dedicado a la literatura, como yo mismo, mucho esfuerzo muy mal pagado. ¿Es solo que no hemos podido evitarlo o hay detrás un niño, un muchacho, alguien que está deslumbrado? ¿Por qué no nos hemos dedicado a lo nuestro, que lo hacemos bien y no nos complica la vida?}

Mirá, hablabas recién de mi mujer. Así como nunca fui infiel en mi relación de pareja, siempre fui infiel en mis relaciones laborales, de alguna manera, y creo que en ese sentido la infidelidad tiene un saborcito atractivo. Pero esto lo descubrí hace poco, hace dos, tres, cuatro años, en mi segundo periodo terapéutico. A mí lo que me gusta es aprender. Disfruto como un niño aprendiendo y voy a morir así, voy a morir aprendiendo. Lo que realmente disfruto es eso, y esto de moverme de disciplina, lo que hace es mantenerme todo el tiempo en un aprendizaje constante, ser todo el tiempo un aprendiz.

\preg{Un niño que no quiere crecer. Una especie de Peter Pan intelectual.}

Si querés, vamos a llamarlo así.

\preg{Que es la única manera de no acabar convertido en un fósil, en el capitán Garfio.}

Yo te diría que es eso, pero no fue algo que me propuse. Es que si no tenés ese apetito de conocer o de aprender, la curiosidad, es imposible. Aprender algo es un esfuerzo enorme. Depende de lo que sea, puede ser mayor o menor, pero siempre es un esfuerzo. Si no tenés algo interno que te está empujando a eso, no hay forma de que lo hagas. Y al mismo tiempo, cuando lo haces, es lo que te mantiene vivo y lo que te mantiene joven.

\preg{Aquí hay una vuelta de tuerca que tiene que ver con lo que decías de las infidelidades. Yo llegué a formular que la física fue siempre mi esposa oficial y la literatura era la otra, y siempre me he sentido un poco pecador. Todo lo que dices explica que a mis años siga apasionado con la física, pero ¿por qué además no me canso de la literatura, de entrevistar a gente que me gusta, de la poesía? Es como si no solo nos comiera la curiosidad, sino que el sentido de lo bello aumentara también con la edad.}

No sé si aumenta o simplemente uno aprende a ver la belleza donde antes no la veía. Yo creo que es la apertura, un poco lo que empezamos hablando al principio. Estar abierto te puede deparar sorpresas interesantes, y yo creo que ser abierto te premia bastante después. La apertura mental, estética, social, humana, todo eso no puede más que depararte alegrías al final.

\preg{Vamos a cosas más concretas. \emph{La entrevista}, la obra de teatro, es en mi opinión una de las cosas mayores que has hecho. Yo te conocí a través de ella, antes que al personaje de carne y hueso, y me pareció algo completamente fuera de lo normal por la audacia que suponía. Quiero saber el cómo, el porqué y cómo la ves desde la distancia.}

Te agradezco la pregunta, porque coincido plenamente contigo. Yo creo que \emph{La entrevista} fue una de las cosas más importantes que hice, si no la que más, y que además me cambió un montón de cosas. Empezó como una propuesta para Passion for Knowledge de 2013. Un año antes, en 2012, estábamos terminando con una experiencia piloto, que era el programa de escritores en residencia, en el que estuvo Eduardo Berti, y a partir de ahí surgió lo del proyecto de \emph{La entrevista}. Surgió, yo te diría, casi como una locura. Desde que surgió el proyecto teníamos catorce meses para poner una obra encima de un escenario, y no de cualquier escenario, del Victoria Eugenia, que no era un teatro de barrio. Y no teníamos nada. Cuando te digo nada es que no teníamos experiencia, no teníamos guion, no teníamos actores, no teníamos director, no teníamos por supuesto vestuario, no teníamos ni compañía. Cuando lo miro retrospectivamente y pienso lo que logramos en catorce meses, me parece increíble, y también el nivel de inconsciencia nuestro de habernos metido en esa aventura. Eso por un lado. Por otro lado fue una experiencia absolutamente increíble en todos los aspectos. La construcción de la obra la hicimos Luisa Etxenike y yo. Luisa ya tenía muchos libros publicados, era una escritora consagrada. Yo ya había escrito algo de ficción, pero era como mi primera experiencia literaria seria, y para ella era su primera experiencia compartiendo coautoría, nunca había escrito a cuatro manos. Luisa y yo ya nos conocíamos desde hacía unos años, pero el punto crucial era encontrar un lenguaje común. Una vez que ese lenguaje común apareció, todo lo demás fue bastante fluido. De hecho la escribimos en este mismo edificio, en el piso de arriba, en una oficina allí al fondo que nos había dado el DIPC. Nos reuníamos dos o tres veces por semana, a veces dos horas, a veces todo el día. Y lo más interesante de todo fue que es realmente una escritura a cuatro manos. No fue que la queríamos a Luisa para que le diera un toque literario a una tontería científica que yo pudiera escribir, ni tampoco ella necesitaba un escenario científico para que unos personajes literarios deambularan por ahí. Realmente era una construcción de a dos.

\preg{Que no es poco, porque implica salirse del propio ego.}

Totalmente. Y yo creo que salió muy bien por algo que me surgió después de haber escrito \emph{La entrevista} y que es para mí la consigna fundamental de Mestizajes, que es qué podemos hacer juntos tú y yo que ninguno de los dos puede hacer por separado. Esa es la idea fundamental, porque si no, no tiene sentido que nos juntemos a hacer nada. Y eso, sin haberlo hecho consciente en ese momento, fue así. Fue tan escrita a cuatro manos que yo hoy por hoy no te puedo decir si esta oración la escribió Luisa o la escribí yo.

\preg{Y sin embargo, para el lector externo hay un solo autor. Los dos autores desaparecen.}

Lo que también me gusta de este proyecto, y en realidad de todos los proyectos de Mestizajes, porque siempre hay por lo menos dos personas involucradas y en general más, es que ayuda a erradicar del imaginario colectivo la idea de la creatividad individual, porque eso no existe. Uno vive en un entorno, en un contexto, y ese contexto te alimenta necesariamente. La creatividad individual es una cosa que definió el Romanticismo en una época. Esto pasa también en la ciencia, lo sabemos perfectamente. Ninguna de las grandes mentes de la ciencia plasmó ideas totalmente ajenas a la época en que trabajó. A mí me parece que hay que desandar un poco esa idea de la creatividad individual, porque es un invento y no existe.

\preg{Conseguís vuestro libreto. ¿Cómo cerráis el resto de los milagros?}

Con una enorme suerte y con carambolas varias. Fuimos trabajando el guion y, cuando teníamos una versión relativamente avanzada, empezamos a tantear a gente que pudiera hacer la dirección, la puesta en escena. Conocimos al que fue después el director, Pablo Viar. Curiosamente, éramos todos de alguna manera aprendices. Pablo había sido asistente de dirección de varias obras, pero nunca había dirigido en solitario. Fue la primera obra que dirigió en solitario. Pero él ya conocía actores, conocía técnicos. Después consiguió a los dos actores, que fueron David Luque y Aitor Mazo, que falleció creo que al año siguiente o al año y medio de estrenar la obra, y a todo el equipo técnico, que también lo trajo él, no sé si del Centro Dramático Nacional o por ahí, que es la parte del espacio sonoro, de la iluminación, del sonido y de todo. Eran dimensiones que yo desconocía. Cuando nosotros terminamos de escribir el texto, lo podemos considerar como una especie de construcción en 2D, y luego la iluminación, el sonido, la dirección escénica lo convirtieron en 3D.

\preg{¿La experiencia fue transformadora para ti?}

Para mí fue increíble. Además, por temas que se tratan en la obra, me cambió la manera de relacionarme. Yo estaba en un momento de mi carrera donde tenía gente por encima y gente por debajo, y me cambió totalmente la manera de relacionarme tanto con los de arriba como con los de abajo. Por ejemplo, el tema de la puerta del despacho. Hasta escribir la obra yo trabajaba con la puerta cerrada, y hay una frase en la obra que dice «y además sabía que nunca cerrabas la puerta de tu despacho con llave». Y yo ahí dije «joder». Es como el lenguaje corporal, pero esto sería un lenguaje arquitectónico, si se quiere. Trabajar con la puerta abierta o con la puerta cerrada, que es decir «aquí no entres». A partir de ahí, salvo raras ocasiones en que tengo que pensar o estar concentrado y la cierro, y mientras no haya ruido en el pasillo, prefiero trabajar con la puerta abierta.

\preg{Dado el éxito de la obra y su repercusión, ¿por qué no ha habido más obras de teatro?}

Yo creo que tiene que ver con todo lo que estuvimos hablando. Lo hicimos bien, pero debo admitir que tuvimos suerte.

\preg{Y uno no quiere tentar la suerte dos veces.}

Exactamente. No somos personas del ámbito del teatro. Luisa y yo, o yo con alguien más, podríamos escribir otra obra de teatro, pero es otra vez empezar de cero, otra vez buscar una compañía. Es un mundo que nos resulta ajeno, en el que en ese momento nos supimos mover, nos ayudó la suerte y salió bien, y ya está. Y aparte volvemos a esto de la infidelidad y de los cambios. Teatro ya hice, vamos a otra cosa. Después de la obra de teatro hice audiovisual, después hicimos la exposición fotográfica, lo de Creativium.

\preg{Oyéndote me doy cuenta de que todo lo que dices de la obra de teatro te lo podría haber dicho yo del proyecto NEXT, que se hace posible gracias a la absoluta ignorancia de dónde se está metiendo uno. Francis Halzen, cuando lo entrevisté, me contaba de AMANDA, luego IceCube, el experimento de neutrinos más maravilloso del universo, que ellos no se habían leído los libros, que decían que la luz en el hielo se propaga como mucho unas pocas decenas de metros. Resultó que en el hielo ultrapuro de la Antártida, que nadie había medido, se propagaba cientos de metros. Da que pensar que en ciencia, como en arte o en literatura, haga falta la combinación de lo insensato y lo afortunado.}

Es que el azar, y eso es una cosa de la que hablo mucho en el libro, fue como un inquilino incómodo en el palacio de la ciencia durante muchos siglos, y es hora de que aprendamos a vivir con él y a aprovecharlo, porque es lo que permite que surjan un montón de cosas. Hago un paréntesis con otra cosa, el tema de los descubrimientos simultáneos, que se creía una curiosidad y cada vez hay más. Hay una página de Wikipedia que se llama «List of multiple discoveries» con un listado enorme, porque los descubrimientos al final tienen que ver con un contexto histórico que te lleva a ellos. Pero me desvié de lo del azar. Volviendo a los escritores, hay un escritor argentino contemporáneo que a mí me gusta mucho, que es Martín Caparrós, y en la novela \emph{Valfierno} hay un párrafo, que cito en el libro, que raya en la filosofía. Él habla mucho del azar, igual que escritores como Paul Auster, que tiene \emph{La música del azar} y muchas otras cosas acerca del azar. ¿Sabés cómo es lo de Paul Auster? Cuando era chico, tenía catorce años, estaba en un campamento de verano y tenían que pasar por debajo de una alambrada. Iban uno detrás de otro, y justo cuando pasa el chico que iba delante de él, cae un rayo y lo mata. Eso le pasó de verdad, ¿cómo no te va a marcar? Está en mucha de la obra de Paul Auster, e incluso Woody Allen, en \emph{Match Point}, trata de cómo el azar te cambia. Y Caparrós habla maravillosamente de la tontería del azar y del accidente, que es cuando el azar se vuelve pomposo. Él dice que el azar está todo el tiempo, pero que no nos damos cuenta, y que de pronto cambia algo radicalmente, te cambia la vida en un segundo, y a eso le llamamos accidente. Yo creo que el azar se debería enseñar en el colegio, junto con la teoría de la probabilidad, aparte de que es la esencia que mueve la evolución.

\preg{Eso requeriría otra entrevista entera. A mí me pasa que no puedo viajar a París sin pensar en la Maga. ¿Y a ti?}

Hace muchos años, por 2009, hicimos un viaje a París y yo me propuse hacer el circuito que se menciona en \emph{Rayuela}. Me fui al canal Saint-Martin, me fui a un montón de lugares y me saqué una foto en el Pont Neuf igual a la que le hizo Sara Facio a Cortázar. De hecho tengo un libro que se llama \emph{El París de Cortázar}, que hace todo un circuito, y creo que si vas a París hay tours organizados. Yo creo que va a ser lo que haga en el próximo viaje. Son indisolubles París y Cortázar. Y hay una historia muy curiosa de algún viaje anterior, de 2006 o 2007. Estábamos en París el último día de vacaciones y le digo a Silvina «me parece recordar, de alguna cosa que leí en algún lado, que Cortázar está enterrado en el cementerio de Montparnasse, ¿me acompañás?». Era una cosa vaga. Aparcamos por ahí cerca y yo pensaba «¿cómo voy a encontrar la tumba de Cortázar acá?», confiado en que era el único lunático que iba a verla. Y entrás al cementerio de Montparnasse y hay un cartel enorme, «busque a su difunto favorito», y están Sartre, Simone de Beauvoir, están todos. Buscás ahí el sector, vas y la encontrás. A partir de ahí empecé un ciclo de lo que yo llamo turismo necrológico. Estuve en la tumba de Borges en Ginebra.

\preg{Yo estuve hace unas semanas con mi hijo.}

No sé si conocés la historia de esa lápida. Yo publiqué un post en mi blog y está toda la historia ahí.

\preg{Mi hijo me preguntó qué eran esas frases en inglés y nórdico antiguos. Le dije lo que sabía, que a Borges le encantaban las lenguas germánicas antiguas, pero ni idea.}

Tiene una historia, en la que Borges, plenamente consciente de lo que estaba haciendo, está labrando su propio mito de alguna manera. Y después estuve en la de Gödel en Princeton, naturalmente buscando la de Einstein, hasta que me enteré de la historia verdadera del cuerpo de Einstein, que había pedido que lo cremaran y lo tiraran al río, no sé cómo se llama el río de Princeton. A falta de Einstein me fui a ver la de Gödel, que no está mal tampoco. Esa es mi historia de turismo necrológico.

\preg{De Cortázar siempre me ha impresionado su internacionalidad. Era argentino hasta que era francés, y francés hasta que era internacional. En «La noche boca arriba» ves a la vez a un chico en moto y a los pueblos precolombinos, en \emph{Rayuela} la bohemia y la intelectualidad, en «La autopista del sur» un contacto humano casi de ciudad pequeñita, y «Continuidad de los parques» tiene algo de cuento de misterio a la inglesa. Tiene tantas aristas que es difícil capturarlo, y sin embargo uno sabe que es un escritor argentino que escribe en español y no un escritor francés.}

Ni siquiera que escribe en español, porque Cortázar escribe en argentino. Yo, por ejemplo, y evitando cualquier comparación con Cortázar, por favor, que quede claro, ya no puedo escribir en argentino. Me cuesta horrores, salvo que un personaje lo necesite. Me acostumbré a conjugar en tú y la conjugación en vos me suena un poco forzada. En una conversación, sobre todo si estoy con argentinos, obviamente es completamente natural, pero cuando lo tengo que escribir me cuesta, salvo que haya personajes que lo requieran. Tengo alguna historia que si el personaje no es argentino ya pierde toda la gracia, porque lo estás identificando por eso y por la forma en que habla. Si te acordás de ese cuento de Cortázar, yo creo que era «La señorita Cora», me parece, en el que juega con el tú y el usted, «usted se tendió a tu lado». Si no hay un personaje argentino no se entiende. De hecho he leído cosas sobre cuando intentan traducir ese cuento, y es una locura. Entonces yo creo que Cortázar siempre escribió en argentino.

\preg{Y siendo en ese sentido más argentino que ningún argentino, sabe ser más internacional que nadie.}

Sí, busca esa universalidad de alguna manera.

\preg{Hay un Cortázar, el del tipo que se está afeitando y es la imagen del espejo la que se corta, con una capacidad para generar tensión, misterio y tragedia que no sé si se había visto antes.}

Yo no sé si la originalidad de Cortázar pasa por la tensión. Yo creo que su gran acierto es lo que él mismo define. Él define lo fantástico no como esas historias de fantasmas, o yéndonos a la ciencia ficción o cosas por el estilo, sino como una grieta en lo cotidiano, una grieta que te introduce en lo fantástico. Muy poca gente dio con esa veta. A mí me gusta mucho la definición. Leí muchos ensayos de él también, y para él lo fantástico está ahí, a medio paso de lo real. Para poner un ejemplo contemporáneo, ¿has visto capítulos de \emph{Black Mirror}? Lo que más pavor te da de \emph{Black Mirror} es que está ahí nomás, que eso puede ocurrir dentro de cuatro o cinco años.

\preg{Se diría que nos hemos olvidado un poco de Cortázar. ¿Es verdad o me lo estoy imaginando?}

Coincido contigo. Me parece que Cortázar fue un poco una singularidad y que no creó escuela, y no es culpa de él. Me parece que no hubo una corriente literaria de lo fantástico en ese sentido. Quizás Calvino, que era de la misma época, pero tampoco. Por el motivo que sea, eso no caló en el mundo literario, o en la sociedad incluso. Ahí habría que preguntarle a algún sociólogo de la literatura qué pasó, porque yo creo que lo que se consume hoy es novela histórica, en el mejor de los casos. No tengo muy claro por qué. Quizás porque se pasó de inquietante, o porque vivimos en un mundo inquietante y entonces queremos anclarnos a ese pasado. Es filosofía de café, pero creo que puede ir por ese lado. Y después está otro tema fundamental, que no es una lectura fácil. Para leer a Cortázar hoy hay que hacer un poco de trabajo.

\preg{Me ha impresionado mucho la idea de que Mestizajes apunte ahora hacia una parte cuantitativa. Te arriesgas a abrir un campo nuevo de la ciencia, o a profundizar en él, y a que en algún momento alguien te riña por haber dejado atrás lo cómodo, los juegos florales, el «voy a invitar aquí a un artista».}

No te creas, porque siempre mantenemos esa interacción con la sociedad y con lo artístico. Lo que hago en \emph{Conexiones}, usar las redes complejas para estudiar fenómenos sociales, literarios o artísticos, da como resultado unas visualizaciones preciosas, visualmente son superatractivas, y cuando las vi dije «esto tiene que ir a una exposición». Y de hecho vamos a hacer una expo, ya está todo confirmado. El 23 de febrero se inaugura en el Museo San Telmo de San Sebastián. A mí me interesa esa exposición, primero, porque es visualmente atractiva y, segundo, porque así como se presentan estos casos más sesudos del arte, la ciencia, la literatura o la filosofía, tú puedes crear otras visualizaciones. Yo lo hice para jugar. Puedes poner cuáles son las relaciones entre Taylor Swift y Nietzsche, o entre Maradona y el papa, o lo que quieras. Vamos a poner todas esas visualizaciones ahí y finalmente va a haber una última estación, que estamos haciendo ahora mismo con una empresa de software que está diseñando una página web, donde la gente va a poder poner dos, tres o cuatro semillas y crear su propio mapa cultural a partir de ellas. Querés poner a Bad Bunny con Madonna, lo ponés ahí y te da el mapa.

\preg{O con Gödel, o lo que sea. Es fantástico.}

A mí siempre me preocupa eso. Siempre busqué, en todas las cosas que hicimos en Mestizajes, que la gente participe.

\preg{Porque la recepción pasiva es «cuéntame el rollo que tú quieras, que yo miro para otro lado».}

Exacto. Siempre evité esa participación pasiva. Cuando había charlas, había un espacio muy generoso para preguntas, discusión y debate. Y aquí, además de participar y ver, la idea es que puedan jugar. Va a haber un código QR y la gente se lo lleva en el móvil. Puedes ir en el autobús creando el mapa de Bad Bunny con Madonna, aprietas un botón y lo mandas por las redes sociales, lo compartes o lo que sea. Y los investigadores van a poder descargarse esa red, que viene definida por una matriz, y trabajar con ella.

\preg{Esta descripción me lleva a la última pregunta, que no estaba preparada. Si en una de tus redes pinchamos Messi y Dios, ¿qué va a pasar?}

Lo podemos probar. Yo puse a Maradona y al papa, que no andan muy lejos. Es curioso, porque me salen dos clústeres muy separados, pero hay tres o cuatro nodos ahí en el medio y son los últimos tres presidentes argentinos. Es increíble. Y cuando pones a Taylor Swift y a Nietzsche también quedan dos clústeres muy separados, y cuando ves cuáles son los elementos del ámbito de la filosofía que están más cerca del universo de Taylor Swift y compañía, aparece el hedonismo, aparece la estética, aparece la alienación social.

\preg{Has dado con algo realmente flipante.}

Flipante. Es realmente flipante.

\preg{Pues creo que este es un buen momento para dejar la entrevista, con la invitación a la gente a que venga a la exposición, se traiga su teléfono y nos conteste cuál es la conexión entre Messi y Dios, que va a ser extremadamente interesante. Muchas gracias.}

Un placer, Juanjo. Un placer muy grande.

\end{document}
